\expandafter\ifx\csname CAPETINCLUDEMODE\endcsname\relax
\documentclass[11pt,a4paper]{article}
% preamble.tex — packages et commandes communes a toutes les fiches CAPET SII IE
\input{glyphtounicode}
\pdfgentounicode=1
\usepackage{cmap}
\usepackage[T1]{fontenc}
\usepackage[utf8]{inputenc}
\usepackage[french]{babel}
% babel[french] active a \begin{document} des raccourcis de ponctuation (:;!?)
% qui cassent la syntaxe to[...] de circuitikz (erreur "+ or - expected") ;
% on les desactive, ce qui ne change que l'espacement fin avant ces signes.
\AtBeginDocument{\shorthandoff{:;!?}}
\usepackage[a4paper, top=1.6cm, bottom=1.6cm, left=1.4cm, right=1.4cm]{geometry}
\setlength{\headheight}{14.5pt}
\usepackage{amsmath, amssymb, amsfonts}
\usepackage{siunitx}
\sisetup{locale=FR, output-decimal-marker={,}, inter-unit-product=\ensuremath{{}\cdot{}},
         group-minimum-digits=4}
% Unites non fournies par siunitx, declarees ici une fois pour toutes (fiches + exercices).
% Les gardes \ifdefined evitent tout conflit avec une redeclaration locale restee dans un
% fichier : ne pas retirer les gardes, ne pas redeclarer ces unites dans les fichiers.
\DeclareSIUnit{\jour}{j}
\ifdefined\var\else\DeclareSIUnit{\var}{var}\fi   % puissance reactive
\ifdefined\vad\else\DeclareSIUnit{\vad}{VAD}\fi   % puissance deformante
\ifdefined\tr\else\DeclareSIUnit{\tr}{tr}\fi      % tour (tr/min)
\usepackage{array, booktabs, multirow, tabularx}
\usepackage{xcolor}
\usepackage[most, breakable]{tcolorbox}
\usepackage{multicol}
\usepackage{enumitem}
\usepackage{mathtools}
\usepackage{lmodern}
% Refonte visuelle 2026 (direction "moderne coloré") : Lato pour le texte,
% maths sans-serif assorties via newtxsf (charge apres amsmath/amssymb).
\usepackage[default]{lato}
\usepackage{newtxsf}
\usepackage{microtype}
\usepackage{fancyhdr}
\usepackage{tikz}
\usetikzlibrary{positioning,arrows.meta}
\usepackage{circuitikz}
\usepackage{lastpage}

% ── Couleurs par domaine (palette modernisée 2026) ────────────────────
\definecolor{colElec}{RGB}{29, 78, 216}        % Électrotechnique — bleu vif
\definecolor{colElecLight}{RGB}{219, 234, 254} % fond bleu clair
\definecolor{colAuto}{RGB}{21, 128, 61}        % Automatique — vert
\definecolor{colMeca}{RGB}{194, 65, 12}        % Mécanique — orange
\definecolor{colTherm}{RGB}{185, 28, 28}       % Thermique — rouge
\definecolor{colQT}{RGB}{232, 244, 233}        % (hérité, plus utilisé par questiontype)
\definecolor{colJury}{RGB}{255, 248, 225}      % fond jury (crème)
\definecolor{colJuryFrame}{RGB}{202, 138, 4}   % bordure/bandeau jury (ambre)
\definecolor{colGrayLight}{RGB}{245, 245, 247} % fond gris neutre
\definecolor{colSeparator}{RGB}{180, 180, 180} % gris séparateur
\definecolor{colExo}{RGB}{8, 145, 178}          % Banque d'exercices — cyan (distinct des fiches)
% ── Encre et gris de la refonte ───────────────────────────────────────
\definecolor{colInk}{RGB}{23, 32, 42}          % texte courant (quasi-noir)
\definecolor{colSlate}{RGB}{83, 98, 118}       % gris ardoise (méta, en-têtes)
\definecolor{colRule}{RGB}{203, 213, 225}      % filets clairs
\definecolor{colJuryBg}{RGB}{255, 248, 225}    % fond crème jury
\definecolor{colPresence}{RGB}{83, 98, 118}    % bloc 5 (présence dans les sujets)
\AtBeginDocument{\color{colInk}}

% ── Style encadré exercice + corrigé (banque d'exercices) ─────────────
\tcbset{
  exoBox/.style={
    breakable,
    enhanced,
    colback=colExo!8!white,
    colframe=colExo,
    colbacktitle=colExo,
    coltitle=white,
    fonttitle=\bfseries\small,
    left=4pt, right=4pt, top=3pt, bottom=3pt,
    before skip=5pt, after skip=5pt,
  }
}

% ── Bascule enonce / corrige (banque d'exercices) ─────────────────────
% \CAPETCORRIGE est defini en ligne de commande (pdflatex "\def\CAPETCORRIGE{1}\input{...}")
% pour compiler la variante corriges d'un meme fichier source ; par defaut (non defini),
% la compilation produit la variante enonces. \CAPETSHOW{enonce}{corrige} choisit la branche.
%
% TROIS MODES (refonte 2026-09-20, manuel de revision) :
%   - aucun drapeau   -> imprime l'enonce            (variante enonces, INCHANGE)
%   - \CAPETCORRIGE   -> imprime le corrige          (variante corriges, INCHANGE)
%   - \CAPETMANUEL    -> imprime l'enonce ET met le corrige de cote pour l'annexe
%
% Le mode manuel stocke chaque corrige dans une macro (\csgdef), sans ecriture
% fichier ni re-tokenisation : le contenu (circuitikz, accents, %) traverse intact.
% L'annexe les rejoue ensuite dans l'ordre via \CAPETANNEXE.
\usepackage{etoolbox}

\newcounter{CAPETexo}          % numero continu d'exercice dans le manuel
\setcounter{CAPETexo}{0}
\newcounter{CAPETannexe}       % compteur de relecture en annexe

% Titre courant de la notion en cours (pose par \CAPETNOTION), sert d'intertitre annexe
\newcommand{\CAPETnotioncode}{}
\newcommand{\CAPETnotiontitre}{}

% \CAPETNOTION{code}{titre} : ouvre une notion dans le manuel (memorise le libelle,
% et marque le rang du premier exercice de la notion pour l'intertitre d'annexe).
\newcommand{\CAPETNOTION}[2]{%
  \renewcommand{\CAPETnotioncode}{#1}%
  \renewcommand{\CAPETnotiontitre}{#2}%
  \csgdef{CAPETnotionstart\the\numexpr\value{CAPETexo}+1\relax}{#1 -- #2}%
}

% ── \degree robuste ───────────────────────────────────────────────────
% Certains fichiers d'exercices utilisent \degree. Ni le preambule ni siunitx
% ne le definissent, et \providecommand local ne suffit pas : sous hyperref,
% \degree existe deja comme commande factice (\ERROR), donc \providecommand
% s'abstient et l'appel echoue a la compilation du manuel (erreur invisible
% en compilation isolee). On le (re)definit donc inconditionnellement ici,
% apres le chargement des packages, pour tout le corpus.
\AtBeginDocument{\providecommand{\degree}{}\renewcommand{\degree}{\ensuremath{^\circ}}}

% Libelle decoratif : les fichiers d'exercices utilisent aussi \CAPETSHOW dans
% leur \fancyhead et leur bloc de titre, pour afficher "(enonces)" ou "(corriges)".
% Ces appels-la ne sont PAS des exercices : en mode manuel ils ne doivent ni
% consommer un numero, ni alimenter l'annexe. On les reconnait a leur premier
% argument, strictement egal au libelle ci-dessous.
\newcommand{\CAPETLIBELLEENONCES}{(\'enonc\'es)}

\ifdefined\CAPETMANUEL
  % --- Mode manuel : enonce en place, corrige differe vers l'annexe ---
  \newcommand{\CAPETSHOW}[2]{%
    \def\CAPETargun{#1}%
    \ifx\CAPETargun\CAPETLIBELLEENONCES
      #1%            appel decoratif : rendu tel quel, hors numerotation
    \else
      \CAPETSHOWEXO{#1}{#2}%
    \fi
  }
  \newcommand{\CAPETSHOWEXO}[2]{%
    \stepcounter{CAPETexo}%
    \csgdef{CAPETcorrige\theCAPETexo}{#2}%
    \edef\CAPETtmpnum{\theCAPETexo}%
    #1%
    \par\vspace{2pt}%
    {\footnotesize\itshape\color{colExo}%
     \hyperref[capet:cor:\CAPETtmpnum]{Corrige de l'exercice \CAPETtmpnum\ $\rightarrow$ p.~\pageref{capet:cor:\CAPETtmpnum}}%
     \label{capet:eno:\CAPETtmpnum}}%
  }
  % \CAPETANNEXE : rejoue tous les corriges captures, dans l'ordre, avec intertitres
  \newcommand{\CAPETANNEXE}{%
    \setcounter{CAPETannexe}{0}%
    \loop\ifnum\value{CAPETannexe}<\value{CAPETexo}%
      \stepcounter{CAPETannexe}%
      \edef\CAPETcurnum{\theCAPETannexe}%
      \ifcsname CAPETnotionstart\CAPETcurnum\endcsname
        \par\medskip
        {\normalsize\bfseries\color{colExo}\csuse{CAPETnotionstart\CAPETcurnum}}%
        \par\nobreak\smallskip
      \fi
      \begin{tcolorbox}[exoBox, breakable,
        title={Corrige de l'exercice \CAPETcurnum\
               \hfill\normalfont\footnotesize
               \hyperref[capet:eno:\CAPETcurnum]{enonce p.~\pageref{capet:eno:\CAPETcurnum}}}]%
      \label{capet:cor:\CAPETcurnum}%
      \csuse{CAPETcorrige\CAPETcurnum}%
      \end{tcolorbox}%
    \repeat
  }
\else
  % --- Modes historiques, comportement strictement inchange ---
  \newcommand{\CAPETSHOW}[2]{\ifdefined\CAPETCORRIGE #2\else #1\fi}
  \newcommand{\CAPETANNEXE}{}
\fi

% ── Difficulté 1 à 3 (étoiles pleines, amssymb) ───────────────────────
\newcommand{\difficulte}[1]{%
  \ifcase#1\relax
  \or $\bigstar$\;niveau 1/3%
  \or $\bigstar\bigstar$\;niveau 2/3%
  \or $\bigstar\bigstar\bigstar$\;niveau 3/3%
  \fi
}

% ── Style encadré extrait de document technique (banque d'exercices) ──
\tcbset{
  docTech/.style={
    breakable,
    enhanced,
    colback=colGrayLight,
    colframe=colSeparator,
    fonttitle=\bfseries\small\itshape,
    coltitle=black,
    colbacktitle=colGrayLight,
    left=4pt, right=4pt, top=3pt, bottom=3pt,
    before skip=4pt, after skip=4pt,
  }
}

% ── Style tcolorbox de section principale ────────────────────────────
% Refonte 2026 : bandeau titre coloré arrondi + fond très pâle, sans cadre.
\tcbset{
  ficheSection/.style={
    breakable, enhanced jigsaw,
    frame hidden, boxrule=0pt, arc=1.8mm,
    colback=#1!4!white,
    colbacktitle=#1,
    coltitle=white,
    fonttitle=\sffamily\bfseries,
    left=7pt, right=7pt, top=5pt, bottom=5pt,
    toptitle=3pt, bottomtitle=3pt,
    before skip=10pt, after skip=8pt,
  }
}

% ── Style encadré question type ──────────────────────────────────────
% Carte blanche, bandeau sombre, liseré très clair.
\tcbset{
  questiontype/.style={
    breakable, enhanced jigsaw,
    colback=white,
    frame hidden, boxrule=0pt, arc=1.8mm,
    colbacktitle=colElec!78!black,
    coltitle=white,
    fonttitle=\sffamily\bfseries\small,
    left=7pt, right=7pt, top=4pt, bottom=4pt,
    toptitle=2.5pt, bottomtitle=2.5pt,
    before skip=8pt, after skip=8pt,
    borderline={0.6pt}{0pt}{colElec!12!white},
  }
}

% ── Style encadré jury ───────────────────────────────────────────────
% Bandeau ambre, fond crème, sans cadre.
\tcbset{
  jury/.style={
    breakable, enhanced jigsaw,
    colback=colJuryBg,
    frame hidden, boxrule=0pt, arc=1.8mm,
    colbacktitle=colJuryFrame,
    coltitle=white,
    fonttitle=\sffamily\bfseries,
    left=7pt, right=7pt, top=4pt, bottom=4pt,
    toptitle=3pt, bottomtitle=3pt,
    before skip=10pt, after skip=8pt,
  }
}

% ── En-tête / pied de page communs (refonte 2026) ─────────────────────
% #1 code fiche, #2 titre court, #3 couleur domaine
\newcommand{\ficheHeaders}[3]{%
  \fancyhead[L]{\small\sffamily{\bfseries\textcolor{#3}{#1}}\hspace{0.7em}\textcolor{colSlate}{#2}}%
  \fancyhead[R]{\small\sffamily\textcolor{colSlate}{CAPET SII IE\, --- \,Session 2026}}%
  \fancyfoot[C]{\small\sffamily\textcolor{colSlate}{\thepage\,/\,\pageref{LastPage}}}%
  \fancyhead[C]{}\fancyfoot[L]{}\fancyfoot[R]{}%
}

% ── Bloc en-tête de fiche (BLOC 0, refonte 2026) ──────────────────────
% #1 code  #2 titre  #3 couleur domaine  #4 domaine  #5 ligne présence  #6 sous-thèmes
\newcommand{\ficheEntete}[6]{%
  \begin{tcolorbox}[enhanced, breakable,
    frame hidden, boxrule=0pt, arc=2.5mm,
    interior style={left color=#3, right color=#3!75!black},
    left=9pt, right=9pt, top=7pt, bottom=7pt,
    before skip=0pt, after skip=6pt]
    {\color{white}\sffamily\bfseries\LARGE #1\quad #2}\\[4pt]
    {\color{#3!22!white}\sffamily\small \MakeUppercase{#4}\quad\textbar\quad #5}\\[1pt]
    {\color{#3!35!white}\sffamily\footnotesize Sous-thèmes : #6}
  \end{tcolorbox}%
}

% ── Mise en page ─────────────────────────────────────────────────────
\pagestyle{fancy}
\fancyhf{}
\renewcommand{\headrulewidth}{0.5pt}
\renewcommand{\footrulewidth}{0pt}
\AtBeginDocument{%
  \renewcommand{\headrule}{{\color{colRule}\hrule width\headwidth height 0.5pt}}%
}

\setlist[itemize]{leftmargin=1.2em, itemsep=1pt, topsep=2pt, parsep=0pt}
\setlist[enumerate]{leftmargin=1.2em, itemsep=3pt, topsep=2pt, parsep=0pt}

\setlength{\parskip}{2pt}
\setlength{\columnsep}{10pt}

\def\CAPETBEGINFILE{\begin{document}}
\def\CAPETENDFILE{\end{document}}
\else
\def\CAPETBEGINFILE{}
\def\CAPETENDFILE{}
\fi

\fancyhead[L]{\textbf{\textcolor{colExo}{T07}} - Exercices - Hacheurs}
\fancyhead[R]{\small\textit{Banque d'exercices \CAPETSHOW{(\'enonc\'es)}{(corrig\'es)}}}
\fancyfoot[C]{\small\thepage\,/\,\pageref{LastPage}}

\CAPETBEGINFILE

\begin{tcolorbox}[
  enhanced, breakable,
  colback=colExo,
  colframe=colExo,
  left=6pt, right=6pt, top=5pt, bottom=5pt,
  before skip=0pt, after skip=4pt
]
  {\color{white}\bfseries\large T07 \quad Hacheurs (1, 2 et 4 quadrants)}\\[2pt]
  {\color{white}\small Banque d'exercices -- 5 exercices \CAPETSHOW{(\'enonc\'es)}{(corrig\'es)}}
\end{tcolorbox}

\begin{tcolorbox}[exoBox,
  title={Exercice 1 -- \difficulte{1} \hfill\normalfont\footnotesize 2025, \'epreuve \'ecrite disciplinaire appliqu\'ee, questions 14 et 19}
]
\CAPETSHOW{
Le convertisseur CV1 adapte la tension d\'elivr\'ee par des panneaux photovolta\"iques $U_{PV}$ \`a la tension du bus de stockage $E_b = \SI{48}{\volt}$ (batterie), dans une cha\^ine de r\'egulation MPPT (recherche du point de puissance maximale). Le sujet lie les deux tensions par une \emph{fonction de modulation} $f_m$ telle que $U_{PV} = f_m\,E_b$.

\begin{center}
\begin{tikzpicture}[scale=1, >=Stealth, node distance=1.4cm]
  \node[draw, align=center, minimum width=2.4cm, minimum height=1cm] (pv) {Panneaux PV\\$U_{PV}$};
  \node[draw, align=center, minimum width=2.6cm, minimum height=1cm, right=1.3cm of pv] (cv1) {CV1\\hacheur, $\alpha$};
  \node[draw, align=center, minimum width=2.6cm, minimum height=1cm, right=1.3cm of cv1] (bus) {Bus DC\\batterie $E_b=\SI{48}{\volt}$};
  \draw[-latex] (pv) -- (cv1);
  \draw[-latex] (cv1) -- (bus);
\end{tikzpicture}
\end{center}

\begin{enumerate}[leftmargin=1.5em]
  \item Trois points de fonctionnement sont relev\'es sur les documents techniques DT6 et DT7 : $P_1$ ($U_1 = \SI{195}{\volt}$), $P_{max}$ (point nominal, $U = \SI{317}{\volt}$) et $P_2$ ($U_2 = \SI{340}{\volt}$). Calculer, pour chacun, la fonction de modulation $f_m = U_{PV}/E_b$.
  \item En d\'eduire la fonction de CV1 : hacheur \'el\'evateur (Boost) ou hacheur d\'evolteur (Buck) ? Justifier.
  \item En admettant la loi du hacheur \'el\'evateur $U_{PV} = E_b/(1-\alpha)$, calculer le rapport cyclique $\alpha$ correspondant \`a chacun des trois points.
\end{enumerate}

}{
\textbf{Corrig\'e.}

\begin{enumerate}[leftmargin=1.5em]
  \item \[ f_{m1} = \dfrac{195}{48} \approx 4{,}06 \qquad f_{m,max} = \dfrac{317}{48} \approx 6{,}60 \qquad f_{m2} = \dfrac{340}{48} \approx 7{,}08 \]
  \item Les trois valeurs sont sup\'erieures \`a 1 : la tension c\^ot\'e panneaux est toujours plus \'elev\'ee que la tension de la batterie $E_b = \SI{48}{\volt}$. La relation $U_{PV} = E_b/(1-\alpha)$ n'admet une solution avec $0 \le \alpha < 1$ que si $U_{PV} \geq E_b$, ce qui est bien le cas ici. Un hacheur d\'evolteur (Buck), \`a l'inverse, ne peut relier que deux tensions telles que $\langle u\rangle = \alpha E \le E$ : il ne pourrait pas satisfaire cette contrainte. CV1 est donc un \textbf{hacheur \'el\'evateur}.
  \item Avec $f_m = 1/(1-\alpha)$, soit $\alpha = 1 - 1/f_m$ :
  \[ \alpha_1 = 1 - \dfrac{1}{4{,}06} \approx 0{,}754 \qquad \alpha_{max} = 1 - \dfrac{1}{6{,}60} \approx 0{,}848 \qquad \alpha_2 = 1 - \dfrac{1}{7{,}08} \approx 0{,}859 \]
  \textbf{Conclusion.} Les trois rapports cycliques restent group\'es entre $0{,}75$ et $0{,}86$ alors que $U_{PV}$ varie de pr\`es de $75\,\%$ (195 \`a \SI{340}{\volt}) : la commande MPPT doit donc r\'esoudre de tr\`es faibles variations de $\alpha$ pour balayer toute la caract\'eristique du champ photovolta\"ique, ce qui fixe la r\'esolution minimale exig\'ee du modulateur de largeur d'impulsion de CV1.
\end{enumerate}
}
\end{tcolorbox}

\begin{tcolorbox}[exoBox,
  title={Exercice 2 -- \difficulte{2} \hfill\normalfont\footnotesize 2025, \'epreuve \'ecrite disciplinaire appliqu\'ee, question 16 (document technique DT5)}
]
\CAPETSHOW{
Le hacheur \'el\'evateur CV1 d'une installation photovolta\"ique en site isol\'e est aliment\'e par un champ de panneaux mont\'es \emph{en s\'erie}, dont la tension nominale vaut $U_{nom} = \SI{317}{\volt}$ et le courant nominal $I_{nom} = \SI{9,47}{\ampere}$. Avant de dimensionner le convertisseur, il faut caract\'eriser sa source.

% Extrait DT5 du sujet E2 disciplinaire IE 2025 : toutes les valeurs ci-dessous sont
% textuellement presentes dans E2_Disciplinaire_IE_2025.md l.599 (caracteristiques d'un panneau)
% et l.587-595 (synoptique : rendement modulateur 90 %, pertes ligne 3 %).
% Seule la mise en forme en tableau est reconstituee (le DT original est une image).
\begin{tcolorbox}[docTech, title={Document technique DT5 -- caract\'eristiques d'un panneau du champ (extrait)}]
\footnotesize
\begin{center}
\renewcommand{\arraystretch}{1.15}
\begin{tabular}{@{}ll@{}}
\toprule
Technologie & monocristallin \\
Rendement du panneau $\eta_{panneau}$ & \SI{18,9}{\percent} \\
Puissance nominale $P_{nom}$ & \SI{375}{\watt} \\
Tension nominale $U_{nom}$ & \SI{39,6}{\volt} \\
Courant nominal $I_{nom}$ & \SI{9,47}{\ampere} \\
Tension \`a vide $U_{co}$ & \SI{48,2}{\volt} \\
Courant de court-circuit $I_{cc}$ & \SI{9,92}{\ampere} \\
Dimensions $L \times l \times e$ & \SI{1980}{\milli\metre} $\times$ \SI{1002}{\milli\metre} $\times$ \SI{40}{\milli\metre} \\
Masse & \SI{22}{\kilo\gram} \\
\midrule
\multicolumn{2}{@{}l}{\textit{Synopt.} : panneaux mont\'es en s\'erie -- modulateur d'\'energie $\eta_{Mod} = \SI{90}{\percent}$} \\
\multicolumn{2}{@{}l}{\hphantom{\textit{Synopt.} : }ligne de distribution : pertes \SI{3}{\percent}} \\
\bottomrule
\end{tabular}
\end{center}
\end{tcolorbox}

\begin{enumerate}[leftmargin=1.5em]
  \item \textbf{Relever} dans le document technique DT5 les trois grandeurs qui permettent de d\'eterminer la composition du champ : la grandeur \'electrique caract\'eristique du mode d'association impos\'e (s\'erie), la puissance unitaire et les dimensions hors-tout d'un panneau.
  \item \textbf{Calculer} le nombre de panneaux du champ, puis en d\'eduire la puissance nominale $P_{nom,champ}$ et la surface totale du champ.
  \item V\'erifier la coh\'erence entre la puissance nominale du champ obtenue \`a la question 2 et le produit $U_{nom}\,I_{nom}$ donn\'e pour le champ complet.
  \item Le hacheur CV1 \'el\`eve la tension du champ vers un bus $E_b$. Sachant que le courant nominal d'un panneau vaut \SI{9,47}{\ampere} et que ce courant traverse \emph{aussi} l'inductance d'entr\'ee du hacheur, indiquer quel calibre de courant permanent il faut au minimum retenir pour l'interrupteur command\'e de CV1, et pourquoi le calibre en \emph{tension} ne peut, lui, pas se d\'eduire de $U_{nom}$.
\end{enumerate}

}{
\textbf{Corrig\'e.}

\begin{enumerate}[leftmargin=1.5em]
  \item Extraction du DT5 : le montage \'etant \emph{en s\'erie}, c'est la \textbf{tension nominale} d'un panneau $U_{nom,1} = \SI{39,6}{\volt}$ qui s'additionne (le courant, lui, est commun \`a tous les panneaux d'une cha\^ine s\'erie). Puissance unitaire : $P_{nom,1} = \SI{375}{\watt}$. Dimensions : $\SI{1,980}{\metre} \times \SI{1,002}{\metre}$.
  \item Nombre de panneaux :
  \[ N = \dfrac{U_{nom,champ}}{U_{nom,1}} = \dfrac{317}{39{,}6} \approx 8{,}0 \;\Longrightarrow\; N = 8 \text{ panneaux en s\'erie} \]
  \[ P_{nom,champ} = 8 \times 375 = \SI{3000}{\watt} = \SI{3}{\kilo\watt} \]
  \[ S = 8 \times 1{,}980 \times 1{,}002 \approx \SI{15,9}{\metre\squared} \approx \SI{16}{\metre\squared} \]
  \item Le champ complet d\'ebite $U_{nom}\,I_{nom} = 317 \times 9{,}47 \approx \SI{3002}{\watt} \approx \SI{3}{\kilo\watt}$ : la valeur co\"incide avec les \SI{3}{\kilo\watt} obtenus par $8 \times \SI{375}{\watt}$, ce qui confirme le montage \emph{tout s\'erie} de 8 panneaux (un montage mixte s\'erie/parall\`ele donnerait la m\^eme puissance mais pas la m\^eme tension \`a \SI{317}{\volt}).
  \item Le montage \'etant enti\`erement s\'erie, le courant du champ est celui d'un seul panneau : $I = \SI{9,47}{\ampere}$ au point nominal, et $I_{cc} = \SI{9,92}{\ampere}$ en court-circuit (irradiance maximale). C'est ce courant qui traverse en permanence l'inductance d'entr\'ee et l'interrupteur command\'e de CV1 pendant sa phase passante : il faut donc un calibre de courant permanent \textbf{sup\'erieur \`a $I_{cc} = \SI{9,92}{\ampere}$}, major\'e de la demi-ondulation de d\'ecoupage et d'une marge de d\'eclassement thermique -- en pratique un composant de l'ordre de \SIrange{15}{20}{\ampere}. En revanche, le calibre en \textbf{tension} ne se d\'eduit pas de $U_{nom} = \SI{317}{\volt}$ : d'une part la tension \`a vide du champ vaut $8 \times 48{,}2 = \SI{385,6}{\volt}$ (par temps froid et fort ensoleillement, panneaux non charg\'es), d'autre part l'interrupteur d'un hacheur supporte \`a l'\'etat bloqu\'e la tension du \emph{c\^ot\'e oppos\'e} du convertisseur, pas celle de sa source.

  \smallskip
  \textbf{Conclusion.} La lecture du DT5 fixe donc les deux contraintes de dimensionnement des semi-conducteurs de CV1 : un calibre en courant impos\'e par $I_{cc} \approx \SI{9,9}{\ampere}$ (grandeur commune \`a tous les panneaux en s\'erie) et un calibre en tension impos\'e par la tension \`a vide du champ, \SI{385,6}{\volt}, et non par sa tension nominale -- omettre cet \'ecart de pr\`es de \SI{70}{\volt} conduirait \`a un claquage de l'interrupteur au premier matin d'hiver ensoleill\'e.
\end{enumerate}
}
\end{tcolorbox}

\begin{tcolorbox}[exoBox,
  title={Exercice 3 -- \difficulte{2} \hfill\normalfont\footnotesize 2023, \'epreuve \'ecrite disciplinaire appliqu\'ee, questions C3 et C4}
]
\CAPETSHOW{
Le convertisseur DC/DC de recharge du robot Bakus transf\`ere l'\'energie d'une batterie LiFePO$_4$ de f.\'e.m. $E_{LiFePO4}$ (r\'esistance n\'egligeable) vers la batterie Li-ion du robot, mod\'elis\'ee par un mod\`ele de Th\'evenin ($E_{Li\text{-}ion}$ en s\'erie avec $R_{Li\text{-}ion}$). Les deux interrupteurs $K_a$ et $K_b$ sont command\'es de fa\c{c}on compl\'ementaire sur la p\'eriode de d\'ecoupage $T$ : $K_a$ ferm\'e et $K_b$ ouvert pour $t \in [0\,;\,\alpha T[$, puis $K_a$ ouvert et $K_b$ ferm\'e pour $t \in [\alpha T\,;\,T[$.

Hypoth\`eses du sujet : interrupteurs parfaits, r\'egime p\'eriodique et conduction continue, $i_{Bakus}$ toujours positif, $E_{Li\text{-}ion}$ constante sur une p\'eriode.

\begin{center}
\begin{circuitikz}[european, scale=0.8]
\draw
(0,0) to[V] (0,4)
(0,4) to[normal open switch, l=$K_a$] (3,4)
(3,0) to[empty diode] (3,4)
(3,0) -- (0,0)
(3,4) to[cute inductor, l=$L$, i^=$i_{Bakus}$] (6,4)
(6,4) to[R, l=$R_{Li\text{-}ion}$] (6,2)
(6,2) to[battery2] (6,0)
(6,0) -- (3,0)
;
\node[left] at (-0.45,2) {$E_{LiFePO4}$};
\node at (0.45,3.2) {\small $+$};
\node at (0.45,0.8) {\small $-$};
\node[left] at (2.8,2) {$K_b$};
\draw[-latex] (3.6,0.35) -- (3.6,3.65);
\node[right] at (3.6,2) {$u_{CS}$};
\node[right] at (6.9,1) {$E_{Li\text{-}ion}$};
\draw (1.5,0) node[ground]{};
\end{circuitikz}
\end{center}

\begin{enumerate}[leftmargin=1.5em]
  \item Pour $t \in [0\,;\,\alpha T[$ : d\'ecrire le sch\'ema \'equivalent du montage (quels \'el\'ements restent dans la maille, quel est l'\'etat de $K_b$ et pourquoi), puis \'etablir l'\'equation diff\'erentielle v\'erifi\'ee par $i_{Bakus}(t)$.
  \item Toujours pour $t \in [0\,;\,\alpha T[$, en supposant d\'esormais $R_{Li\text{-}ion}$ n\'egligeable, d\'eterminer l'expression temporelle de $i_{Bakus}(t)$ en fonction de $E_{LiFePO4}$, $E_{Li\text{-}ion}$, $L$, $t$ et de la valeur initiale $I_{min}$.
  \item Reprendre les deux questions pr\'ec\'edentes pour $t \in [\alpha T\,;\,T[$ : \'equation diff\'erentielle, puis expression temporelle en fonction du temps r\'eduit $t' = t - \alpha T$ et de la valeur $I_{max}$ atteinte en $t = \alpha T$.
  \item En \'ecrivant que le r\'egime est p\'eriodique ($i_{Bakus}(T) = i_{Bakus}(0)$), retrouver la relation entre $\langle u_{Li\text{-}ion}\rangle$, $\alpha$ et $E_{LiFePO4}$ \emph{sans} passer par le trac\'e de $u_{CS}$.
\end{enumerate}

}{
\textbf{Corrig\'e.}

\begin{enumerate}[leftmargin=1.5em]
  \item \textbf{Phase 1 ($K_a$ ferm\'e).} $K_a$ relie directement la source $E_{LiFePO4}$ \`a l'inductance. La tension au point milieu vaut alors $u_{CS} = E_{LiFePO4} > 0$ : la diode $K_b$, dont la cathode est port\'ee \`a ce potentiel et l'anode \`a la masse, est \textbf{bloqu\'ee}. La maille est donc : source $\to$ $K_a$ $\to$ $L$ $\to$ $R_{Li\text{-}ion}$ $\to$ $E_{Li\text{-}ion}$ $\to$ masse. La loi des mailles donne
  \[ L\,\dfrac{\mathrm{d}i_{Bakus}}{\mathrm{d}t} + R_{Li\text{-}ion}\,i_{Bakus} = E_{LiFePO4} - E_{Li\text{-}ion} \]
  \item Avec $R_{Li\text{-}ion}$ n\'eglig\'ee, l'\'equation se r\'eduit \`a $L\,\mathrm{d}i_{Bakus}/\mathrm{d}t = E_{LiFePO4} - E_{Li\text{-}ion}$, second membre constant : le courant cro\^it \textbf{lin\'eairement}.
  \[ i_{Bakus}(t) = \dfrac{E_{LiFePO4} - E_{Li\text{-}ion}}{L}\,t + I_{min} \qquad t \in [0\,;\,\alpha T[ \]
  La pente est positive car $E_{LiFePO4} = \SI{500}{\volt} \gg E_{Li\text{-}ion}$ : c'est la phase de \emph{magn\'etisation} de l'inductance, celle qui charge la batterie du robot.
  \item \textbf{Phase 2 ($K_a$ ouvert, $K_b$ ferm\'e -- roue libre).} L'inductance impose la continuit\'e du courant : celui-ci ne peut plus venir de la source, il se referme par $K_b$. La maille devient $L$ $\to$ $R_{Li\text{-}ion}$ $\to$ $E_{Li\text{-}ion}$ $\to$ $K_b$, et $u_{CS} = 0$ :
  \[ L\,\dfrac{\mathrm{d}i_{Bakus}}{\mathrm{d}t} + R_{Li\text{-}ion}\,i_{Bakus} = -\,E_{Li\text{-}ion} \]
  puis, $R_{Li\text{-}ion}$ n\'eglig\'ee et avec $t' = t - \alpha T$ :
  \[ i_{Bakus}(t') = -\,\dfrac{E_{Li\text{-}ion}}{L}\,t' + I_{max} \qquad t' \in [0\,;\,(1-\alpha)T[ \]
  Le courant d\'ecro\^it lin\'eairement : l'inductance se d\'emagn\'etise en continuant d'alimenter la batterie.
  \item \textbf{P\'eriodicit\'e.} En r\'egime \'etabli, la croissance de la phase 1 compense exactement la d\'ecroissance de la phase 2 :
  \[ \dfrac{E_{LiFePO4} - E_{Li\text{-}ion}}{L}\,\alpha T \;=\; \dfrac{E_{Li\text{-}ion}}{L}\,(1-\alpha)T \]
  \[ \Longrightarrow\; \alpha E_{LiFePO4} - \alpha E_{Li\text{-}ion} = E_{Li\text{-}ion} - \alpha E_{Li\text{-}ion} \;\Longrightarrow\; E_{Li\text{-}ion} = \alpha\,E_{LiFePO4} \]
  Or, $R_{Li\text{-}ion}$ \'etant n\'eglig\'ee, $\langle u_{Li\text{-}ion}\rangle = E_{Li\text{-}ion}$, d'o\`u $\langle u_{Li\text{-}ion}\rangle = \alpha\,E_{LiFePO4}$.

  \smallskip
  \textbf{Conclusion.} La loi du d\'evolteur n'est donc pas un r\'esultat \`a m\'emoriser mais la cons\'equence directe de la condition de p\'eriodicit\'e $\Delta i_{mont\acute{e}e} = \Delta i_{descente}$ : c'est l'\'egalit\'e des aires sous $v_L$ sur une p\'eriode. Cette formulation temporelle est la seule qui permette d'encha\^iner ensuite sur l'ondulation $\Delta i_{Bakus} = (E_{LiFePO4} - E_{Li\text{-}ion})\alpha T / L$ et donc sur le dimensionnement de $L$ impos\'e par l'exigence de pr\'eservation de la batterie Li-ion (\SI{2}{\ampere} cr\^ete \`a cr\^ete au maximum).
\end{enumerate}
}
\end{tcolorbox}

\begin{tcolorbox}[exoBox,
  title={Exercice 4 -- \difficulte{3} \hfill\normalfont\footnotesize 2025, \'epreuve \'ecrite disciplinaire appliqu\'ee, questions 29, 31 et 32}
]
\CAPETSHOW{
Le convertisseur CV3-1 d'une installation en site isol\'e \'el\`eve la tension batterie $E_b$ vers le bus continu $V_{DC}$. Sa structure r\'eelle (document technique DT9) comporte \textbf{deux branches entrelac\'ees} : une premi\`ere cellule ($Q_1$, $D_1$, inductance $L_1$) et une seconde cellule identique ($Q_3$, $D_4$, inductance $L_2$), la commande de la seconde \'etant d\'ecal\'ee de $T/2$ par rapport \`a celle de la premi\`ere. Le rapport cyclique $\alpha$ est commun aux deux branches, tous les interrupteurs sont parfaits et le rendement de CV3-1 est unitaire. Le courant de sortie de l'\'etage est not\'e $i_S$, le courant de batterie $i_b$ (valeur moyenne $I_b$).

On admet pour la premi\`ere branche les r\'esultats \'etablis pr\'ec\'edemment :
\[ i_{L1}(t) = \dfrac{E_b}{L_1}\,t + I_{min}\ \ \text{pour } t \in\, ]0\,;\,\alpha T]
\qquad
i_{L1}(t) = \dfrac{E_b - V_{DC}}{L_1}\,(t-\alpha T) + I_{max}\ \ \text{pour } t \in\, ]\alpha T\,;\,T] \]

\begin{center}
\begin{tikzpicture}[scale=0.9, >=Stealth, font=\footnotesize]
  \draw[->] (0,0) -- (8.6,0) node[right] {$t$};
  \draw[->] (0,-0.2) -- (0,2.3) node[above] {$i_{L1}$, $i_{L2}$};
  \draw[thick, colExo] (0,0.5) -- (2.4,1.8) -- (4,0.5) -- (6.4,1.8) -- (8,0.5);
  \draw[thick, dashed, colElec] (0,1.15) -- (0.4,0.5) -- (2.8,1.8) -- (4.4,0.5) -- (6.8,1.8) -- (8,1.35);
  \draw[dotted] (2.4,0) -- (2.4,1.8);
  \draw[dotted] (4,0) -- (4,0.5);
  \node[below] at (2.4,0) {$\alpha T$};
  \node[below] at (4,0) {$T$};
  \node[below] at (8,0) {$2T$};
  \node[right, colExo] at (8.1,0.5) {$i_{L1}$};
  \node[right, colElec] at (8.1,1.35) {$i_{L2}$};
  \draw[<->] (4,0.5) -- (4.4,0.5) node[midway, below=1pt] {\scriptsize $T/2$};
\end{tikzpicture}
\end{center}

\begin{enumerate}[leftmargin=1.5em]
  \item \'Etablir, sur le mod\`ele de $i_{L1}$, les deux expressions temporelles de $i_{L2}(t)$ : pour $t \in\, ]\tfrac{T}{2}\,;\,\tfrac{T}{2}+\alpha T]$ puis pour $t \in\, ]\tfrac{T}{2}+\alpha T\,;\,\tfrac{3T}{2}]$. Pr\'eciser \`a chaque fois l'\'etat de $Q_3$ et de $D_4$.
  \item Le courant $i_S$ n'est aliment\'e par une branche que pendant la fraction de p\'eriode o\`u la diode de cette branche conduit, et chaque branche transporte la moiti\'e du courant de batterie. En d\'eduire l'expression de la valeur moyenne $\langle i_S\rangle$ en fonction de $\alpha$ et de $I_b$.
  \item En \'ecrivant le bilan de puissance du convertisseur suppos\'e parfait, exprimer $V_{DC}$ en fonction de $\alpha$ et de $E_b$. Conclure sur la nature du convertisseur.
  \item Application num\'erique : $E_b = \SI{48}{\volt}$ et le bus doit \^etre r\'egul\'e \`a $V_{DC} = \SI{230}{\volt}$. Calculer $\alpha$, puis le courant moyen $I_b$ appel\'e \`a la batterie si l'\'etage doit fournir $\langle i_S\rangle = \SI{13}{\ampere}$.
\end{enumerate}

}{
\textbf{Corrig\'e.}

\begin{enumerate}[leftmargin=1.5em]
  \item La seconde branche reproduit le comportement de la premi\`ere, d\'ecal\'e de $T/2$.

  \textbf{Pour $t \in\, ]\tfrac{T}{2}\,;\,\tfrac{T}{2}+\alpha T]$} : $Q_3$ est \textbf{passant}, $D_4$ est \textbf{bloqu\'ee}. L'interrupteur \'etant parfait, $v_{Q3} = 0$, donc $v_{L2} = E_b$ :
  \[ i_{L2}(t) = \dfrac{E_b}{L_2}\left(t - \dfrac{T}{2}\right) + I_{min} \]
  \textbf{Pour $t \in\, ]\tfrac{T}{2}+\alpha T\,;\,\tfrac{3T}{2}]$} : $Q_3$ est \textbf{bloqu\'e}, $D_4$ est \textbf{passante}. La diode \'etant parfaite, $v_{D4} = 0$, donc $v_{L2} = E_b - V_{DC} < 0$ :
  \[ i_{L2}(t) = \dfrac{E_b - V_{DC}}{L_2}\left(t - \alpha T - \dfrac{T}{2}\right) + I_{max} \]
  \item Chaque branche ne d\'ebite vers la sortie que pendant sa phase de roue libre, de dur\'ee $(1-\alpha)T$, et le courant qu'elle transporte alors vaut $I_b/2$ (partage \'egal entre les deux branches identiques). Sur une p\'eriode $T$, les deux branches contribuent :
  \[ \langle i_S\rangle = 2 \cdot \dfrac{1}{T}\,(T - \alpha T)\cdot\dfrac{I_b}{2} = I_b\,(1-\alpha) \]
  \item Convertisseur parfait, donc puissance d'entr\'ee = puissance de sortie :
  \[ P_E = E_b\,I_b = P_S = V_{DC}\,\langle i_S\rangle = V_{DC}\,I_b\,(1-\alpha) \]
  En simplifiant par $I_b \neq 0$ :
  \[ \boxed{V_{DC} = \dfrac{E_b}{1-\alpha}} \]
  Comme $0 \le \alpha < 1$, on a toujours $V_{DC} \geq E_b$ : CV3-1 est bien un \textbf{hacheur \'el\'evateur} (Boost), ici \`a deux cellules entrelac\'ees. Il est remarquable que le bilan \'energ\'etique redonne exactement la loi obtenue par la condition $\langle v_L\rangle = 0$ : les deux voies de d\'emonstration sont \'equivalentes, et le jury accepte l'une comme l'autre \`a condition qu'elle soit men\'ee jusqu'au bout.

  \smallskip
  \textit{Point de vigilance jury (rapport de jury CAPET 2025, commentaires du jury, analyse globale des r\'esultats) : \guillemotleft\ De nombreux candidats ne maitrisent pas les lois basiques de l'\'electricit\'e (la loi des n\oe uds, la loi des mailles, les calculs de puissance ou encore la loi d'Ohm) \guillemotright.}
  \item \[ 1-\alpha = \dfrac{E_b}{V_{DC}} = \dfrac{48}{230} \approx 0{,}2087 \;\Longrightarrow\; \alpha \approx 0{,}791 \]
  \[ I_b = \dfrac{\langle i_S\rangle}{1-\alpha} = \dfrac{13}{0{,}2087} \approx \SI{62,3}{\ampere} \]
  \textbf{Conclusion.} Pour \SI{13}{\ampere} livr\'es au bus continu, la batterie \SI{48}{\volt} doit fournir plus de \SI{62}{\ampere} : c'est la contrepartie in\'evitable de l'\'el\'evation de tension \`a puissance constante ($\approx \SI{3}{\kilo\watt}$). C'est pr\'ecis\'ement ce courant \'elev\'e qui justifie l'entrelacement des deux branches : chacune n'en transporte que la moiti\'e, soit environ \SI{31}{\ampere}, ce qui divise par deux le calibre des semi-conducteurs et par quatre les pertes Joule par branche, tout en doublant la fr\'equence apparente de l'ondulation vue par le condensateur de bus.
\end{enumerate}
}
\end{tcolorbox}

\begin{tcolorbox}[exoBox,
  title={Exercice 5 -- \difficulte{3} \hfill\normalfont\footnotesize exercice in\'edit}
]
\CAPETSHOW{
Un hacheur en pont en H (4 quadrants), aliment\'e par un bus continu $E = \SI{300}{\volt}$, alimente une charge r\'esistive-inductive $R = \SI{8}{\ohm}$, $L$ (frein \`a courant de Foucault). La commande est bipolaire sym\'etrique : les interrupteurs diagonaux ($T_1,T_4$) et ($T_2,T_3$) sont command\'es en compl\'ementaire, avec un rapport cyclique $\alpha$ pour la diagonale $(T_1,T_4)$. Chaque interrupteur command\'e est associ\'e \`a une diode en antiparall\`ele ($D_1$ \`a $D_4$), indispensable au passage du courant inductif lors des commutations.

\begin{center}
\begin{circuitikz}[european, scale=0.85]
\draw (-2.4,4) -- (7.4,4);
\draw (-2.4,0) -- (7.4,0);
\node[above] at (3,4.05) {$+E$};
\node[below] at (3,-0.05) {$0$};
\draw (-2.4,0) to[V] (-2.4,4);
\node[left] at (-3.15,2) {$E$};
\node at (-2.68,3.4) {\small $+$};
\node at (-2.68,0.6) {\small $-$};
\draw (0,4) to[normal open switch] (0,2);
\draw (0,2) to[normal open switch] (0,0);
\node[right] at (0.2,3) {$T_1$};
\node[right] at (0.2,1) {$T_2$};
\draw (-1.2,2) -- (0,2);
\draw (-1.2,2) to[empty diode] (-1.2,4);
\draw (-1.2,0) to[empty diode] (-1.2,2);
\node[left] at (-1.55,3) {$D_1$};
\node[left] at (-1.55,1) {$D_2$};
\draw (6,4) to[normal open switch] (6,2);
\draw (6,2) to[normal open switch] (6,0);
\node[left] at (5.8,3) {$T_3$};
\node[left] at (5.8,1) {$T_4$};
\draw (6,2) -- (7.2,2);
\draw (7.2,2) to[empty diode] (7.2,4);
\draw (7.2,0) to[empty diode] (7.2,2);
\node[right] at (7.55,3) {$D_3$};
\node[right] at (7.55,1) {$D_4$};
\draw (0,2) to[R, l=$R$] (3,2) to[cute inductor, l=$L$] (6,2);
\node[circ] at (0,2) {};
\node[circ] at (6,2) {};
\node[circ] at (-1.2,2) {};
\node[circ] at (7.2,2) {};
\node at (-0.45,1.6) {$A$};
\node at (6.45,1.6) {$B$};
\end{circuitikz}
\end{center}

\begin{enumerate}[leftmargin=1.5em]
  \item Rappeler l'expression de la tension moyenne aux bornes de la charge $\langle u_{AB}\rangle$ en fonction de $E$ et $\alpha$, pour une commande bipolaire sym\'etrique.
  \item En r\'egime permanent, justifier que $\langle u_{AB}\rangle = R\langle i\rangle$ (m\^eme argument qu'avec un hacheur simple quadrant).
  \item Calculer le rapport cyclique $\alpha$ n\'ecessaire pour obtenir un courant moyen $\langle i\rangle = \SI{5}{\ampere}$ dans le frein.
  \item Une commande unipolaire (au lieu de bipolaire) permettrait de doubler la fr\'equence apparente de l'ondulation de courant sans changer la fr\'equence de d\'ecoupage des interrupteurs. Quel est l'int\'er\^et pratique de cette propri\'et\'e pour le dimensionnement de $L$ ?
\end{enumerate}

}{
\textbf{Corrig\'e.}

\begin{enumerate}[leftmargin=1.5em]
  \item Commande bipolaire sym\'etrique (pont en H, 4 quadrants) :
  \[ \langle u_{AB}\rangle = (2\alpha - 1)\,E \]
  qui varie contin\^ument de $-E$ (pour $\alpha=0$) \`a $+E$ (pour $\alpha=1$), en passant par $0$ pour $\alpha = 0{,}5$.
  \item En r\'egime p\'eriodique \'etabli, le courant $i(t)$ dans $L$ v\'erifie $i(T)=i(0)$, donc $\langle L\,\mathrm{d}i/\mathrm{d}t\rangle = 0$ (m\^eme raisonnement que pour un hacheur simple quadrant). En moyennant la loi des mailles $u_{AB} = Ri + L\,\mathrm{d}i/\mathrm{d}t$ sur une p\'eriode, il reste $\langle u_{AB}\rangle = R\langle i\rangle$.
  \item $\langle u_{AB}\rangle = R\langle i\rangle = 8\times5 = \SI{40}{\volt}$, donc :
  \[ 40 = (2\alpha-1)\times300 \;\Longrightarrow\; \alpha = \dfrac{40/300+1}{2} = \dfrac{1{,}133}{2} \approx 0{,}567 \]
  \textbf{Conclusion.} Le point de fonctionnement demand\'e est obtenu tout pr\`es de $\alpha = 0{,}5$, c'est-\`a-dire de la tension moyenne nulle : en commande bipolaire, la charge voit pourtant \`a chaque instant $\pm\SI{300}{\volt}$, donc une ondulation de courant maximale alors que le courant moyen n'est que de \SI{5}{\ampere}. C'est exactement le domaine o\`u la commande bipolaire est la plus p\'enalisante pour le dimensionnement du filtre, ce qui motive la question suivante.
  \item Une ondulation de courant \`a fr\'equence apparente doubl\'ee (\`a fr\'equence de d\'ecoupage $f$ inchang\'ee) permet, pour une m\^eme ondulation maximale tol\'er\'ee $\Delta i_{max}$, de r\'eduire l'inductance $L$ n\'ecessaire (l'ondulation \'etant proportionnelle \`a $1/(Lf_{apparente})$) : c'est le m\^eme b\'en\'efice, en substance, que l'entrelacement de plusieurs bras (exercice 4), obtenu ici sans mat\'eriel suppl\'ementaire, uniquement par le choix de la strat\'egie de commande des 4 interrupteurs.
\end{enumerate}
}
\end{tcolorbox}

\CAPETENDFILE
